\documentclass[12pt]{article}
\usepackage[T1]{fontenc}
\usepackage[utf8]{inputenc}
\usepackage{lmodern}
\usepackage{textcomp}
\usepackage{enumitem}
\usepackage{geometry}
\geometry{margin=1in}
\begin{document}
\begin{center}
Answer Key - Economics - K-12 Level Assessment 19 \
\small (Answers are ordered exactly as in the question paper.)
\end{center}

\section*{Multiple Choice}
\begin{enumerate}[label=\arabic*.]
\item \textbf{Answer:} B
\\ \textit{Options:} A) inflation had already impacted the economy before the fiscal stimulus.; B) the economy initially had some unemployed resources.; C) aggregate supply decreased.; D) aggregate demand is steeply sloped.
\\ \textit{Note:} The correct answer is based on the concept that an expansionary fiscal policy can lead to increased real output by utilizing previously unemployed resources, causing only a small rise in the price level when there is slack in the economy.
\item \textbf{Answer:} C
\\ \textit{Options:} A) increase interest rates and throw the economy into a recession.; B) increase interest rates and depreciate the nation's currency.; C) decrease interest rates and risk an inflationary period.; D) decrease interest rates and throw the economy into a recession.
\\ \textit{Note:} Paying down national debt reduces the supply of government bonds in the market, which can lead to lower interest rates, and when combined with increased government spending from a budget surplus, may fuel inflationary pressures in the economy.
\item \textbf{Answer:} D
\\ \textit{Options:} A) the price of a good divided by its marginal utility.; B) the marginal utility of the good divided by its price.; C) the total utility of the good.; D) the difference between the consumer's value and the market price.
\\ \textit{Note:} Consumer surplus represents the benefit consumers receive when they pay less for a product than the maximum price they are willing to pay, highlighting the difference between their perceived value and the market price.
\item \textbf{Answer:} A
\\ \textit{Options:} A) a monopsony buys from a monopoly; B) a monopoly sells to two different types of consumers; C) a monopoly buys from a monopsony; D) a monopolist sells two different types of goods
\\ \textit{Note:} A bilateral monopoly occurs when a single buyer (monopsony) and a single seller (monopoly) interact in the market, resulting in unique bargaining dynamics that differ from typical competitive or monopolistic scenarios.
\item \textbf{Answer:} B
\\ \textit{Options:} A) nominal GDP falls.; B) the nominal interest rate falls.; C) bond prices fall.; D) the supply of money falls.
\\ \textit{Note:} Households demand more money as an asset when the nominal interest rate falls because lower interest rates reduce the opportunity cost of holding cash, making it more attractive to keep money on hand rather than invest it elsewhere.
\end{enumerate}

\section*{True / False}
\begin{enumerate}[label=\arabic*.]
\setcounter{enumi}{5}
\item \textbf{Answer:} False
\\ \textit{Options:} A) True; B) False
\\ \textit{Note:} Sales tax is considered a regressive tax because it takes a larger percentage of income from low-income individuals than from high-income individuals, as everyone pays the same rate regardless of their income level.
\item \textbf{Answer:} True
\\ \textit{Options:} A) True; B) False
\\ \textit{Note:} Land encompasses all natural resources, such as soil, water, minerals, and forests, that are utilized in the creation of goods and services, making the statement true.
\item \textbf{Answer:} True
\\ \textit{Options:} A) True; B) False
\\ \textit{Note:} Milton Friedman is a prominent advocate of Monetarism, which asserts that the money supply is a key factor influencing economic activity and that government should manage it to stabilize the economy.
\item \textbf{Answer:} False
\\ \textit{Options:} A) True; B) False
\\ \textit{Note:} The stock market primarily facilitates the buying and selling of ownership in companies, while the main source of loanable funds in the financial system comes from savings and deposits in banks and other financial institutions.
\item \textbf{Answer:} True
\\ \textit{Options:} A) True; B) False
\\ \textit{Note:} The statement is true because the substitution effect describes how consumers tend to replace more expensive items with cheaper alternatives when the price of a product decreases, leading them to buy more of the cheaper product.
\end{enumerate}

\section*{Long Form Response}
\begin{enumerate}[label=\arabic*.]
\setcounter{enumi}{10}
\item \textbf{Answer:} A contractionary monetary policy is designed to reduce the money supply in an economy. This policy typically leads to an increase in the discount rate and the nominal interest rate, making borrowing more expensive for banks and consumers. As a result, aggregate demand decreases because higher interest rates discourage spending and investment. The potential implications for the economy include slower economic growth, reduced consumer spending, and lower inflation rates, which can help stabilize an overheated economy but may also lead to higher unemployment in the short term.
\\ \textit{Note:} This answer is correct because it accurately describes how contractionary monetary policy raises the discount and nominal interest rates, leading to decreased borrowing, reduced aggregate demand, and various economic implications such as slower growth and lower inflation.
\item \textbf{Answer:} The downward-sloping nature of the demand curve in economics can be largely explained by substitution effects and income effects. The substitution effect occurs when consumers switch from one product to another when the price of the first product rises, leading to a decrease in quantity demanded. For example, if the price of apples increases, people may choose to buy oranges instead, which are now relatively cheaper. The income effect, on the other hand, happens when a change in price affects consumers' purchasing power; when prices fall, consumers feel richer and may buy more of a good. For instance, if the price of chocolate bars decreases, consumers can buy more bars without spending more money, increasing the quantity demanded. Together, these effects illustrate why demand typically decreases as prices rise, creating a downward slope on the demand curve.
\\ \textit{Note:} correct because it accurately describes the substitution effect, where consumers shift their preferences to relatively cheaper alternatives, and the income effect, where price changes influence consumers' purchasing power, both of which contribute to the overall decrease in quantity demanded as prices rise, thereby explaining the downward slope of the demand curve.
\end{enumerate}

\vfill
\noindent\textit{End of Answer Key}
\end{document}